% --------------------------------------------------------
\chapter{Квантові флуктуації і збурення}
\label{sec:quantum}
%% --------------------------------------------------------

У §\,\ref{sec:inflaton} інфлатон розглядався як класичне однорідне
поле $\bar{\phi}(t)$, а у §\,\ref{sec:metric} ми побудували математичний апарат
для опису збурень метрики. Тепер постає фундаментальне питання про
їхнє походження.

Будь-яке квантове
поле неминуче флуктуює~--- навіть у вакуумному стані. Ці нульові коливання
можна представити як сукупність хвиль (мод Фур'є) найрізноманітніших масштабів~---
від субмікроскопічних до космологічних, де кожна мода характеризується своїм
хвильовим числом $k$ та відповідною фізичною довжиною хвилі $\lambda = 2\pi a(t)/k$.

Під час інфляції ці нульові коливання існують скрізь, зокрема й у нашій
майбутній області простору. Але стрімке інфляційне розширення розтягує
фізичну довжину хвилі кожної моди ($\lambda \propto a(t)$). Щойно для моди
з певним $k$ виконується умова $\lambda \approx cH^{-1}$, вона перетинає горизонт
Хаббла. Протилежні фази цієї хвилі втрачають причинний контакт: вони більше не «взаємодіють» між
собою, і часові осциляції припиняються. Амплітуда хвиль класично
заморожується на тому значенні, яке мода мала в момент виходу з-під горизонту.

Проте горизонт Хаббла не залишається фіксованим назавжди. Після закінчення
інфляції та переходу до радіаційно-домінованої епохи закон розширення Всесвіту
кардинально змінюється. Якщо під час інфляції масштабний фактор зростав
прискорено ($\ddot{a} > 0$), то у гарячу епоху, коли домінує релятивістська
плазма, він сповільнюється як $a(t) \propto t^{1/2}$. Відповідно, параметр
Хаббла спадає як $H = \dot{a}/a = 1/(2t)$.

Це призводить до того, що сопутний горизонт Хаббла, який під час інфляції
стрімко стискався, тепер починає розширюватися:
\begin{equation}
\frac{d}{dt} \left( \frac{c}{aH} \right) > 0.
\label{eq:horizon_growth}
\end{equation}
На мові фізичних відстаней це означає, що реальний радіус горизонту
$R_H = c/H \propto t$ зростає лінійно, випереджаючи ріст довжини хвилі
збурень $\lambda \propto t^{1/2}$.


Внаслідок цієї нерівності раніше заморожені моди одна за одною знову
повертаються під горизонт~--- і саме тут вони «розморожуються», знову
починаючи осцилювати з тією самою амплітудою, яку отримали під час інфляції.
Ця амплітуда стає початковою умовою для коливань гарячої фотон-баріонної
плазми й зрештою відбивається в температурній анізотропії реліктового
випромінювання (CMB) та великомасштабній структурі Всесвіту.


Для того щоб кількісно описати цей складний процес і розрахувати конкретні спостережувані
наслідки, реальні збурення розкладають за типами симетрії простору. Оскільки на ранніх
етапах амплітуди цих коливань є надзвичайно малими ($\delta\rho/\rho \sim 10^{-5}$), еволюцію
різних компонент можна розглядати ізольовано в межах теорії збурень першого порядку.

У лінійному наближенні квантові збурення поділяються на два незалежні класи:
\begin{enumerate}
  \item \textbf{Флуктуації інфлатона} $\delta\phi(\eta, \mathbf{x})$~---
        породжують первинні скалярні збурення кривини, що стають
        «насінням» для формування галактик та їхніх скупчень.
  \item \textbf{Флуктуації метрики} $h_{ij}(\eta, \mathbf{x})$~---
        тензорні збурення простору-часу, які інтерпретуються як первинні
        гравітаційні хвилі.
\end{enumerate}
Ці два класи еволюціонують незалежно і залишають кардинально різні відбитки
у поляризації реліктового випромінювання (так звані $E$- та $B$-моди, які ми детально розглянемо у §\,\ref{sec:pol}).


\textit{Зауваження про квантово-класичний перехід.}
Після виходу за горизонт квантові моди «заморожуються» і
поводяться як класичні стохастичні поля~--- саме їх ми й
спостерігаємо у вигляді температурних плям CMB. Практичним
механізмом, що робить цей перехід ефективним, є \emph{стиснення
квантового стану} (squeezing): у просторі де-Сіттера кожна
суперхаббловська мода еволюціонує у сильно стиснутий когерентний
стан, у якого дисперсія однієї з квадратурних компонент
(«амплітудної») зростає експоненціально як $e^{2Ht}$, тоді як
спряжена («фазова») компонента пригнічується як $e^{-2Ht}$.
У такому граничному стані квантові некомутаційні поправки
$[\hat{\mathcal{R}}_k, \hat{\mathcal{R}}_k'] \sim \hbar$ стають
нехтовно малими порівняно з класичною амплітудою, і мода
практично невідрізнювана від класичного стохастичного поля~---
незалежно від того, чи відбулася повноцінна декогеренція
\cite{polarski_starobinsky_1996, kiefer_polarski_1998}.
Точний механізм цього переходу (вибір базису вказівників,
роль довкілля) залишається відкритим концептуальним питанням
і виходить за рамки цього конспекту. Проте це питання
концептуальне, а не вимірювальне: гауссова статистика заморожених
мод, яку передбачає стандартний підхід, прекрасно узгоджується зі
спостереженнями CMB~--- зокрема з відсутністю первинної
негауссовості за даними Planck~2018 \cite{planck_2018}.
Для практичних розрахунків достатньо того, що амплітуди
заморожених мод підпорядковуються гауссовій статистиці з
дисперсією~\eqref{eq:field_power_spec}.

%% --------------------------------------------------------
\section{Скалярні флуктуації у просторі де-Сіттера}
%% --------------------------------------------------------

Рівняння~\eqref{eq:klein_gordon} описує фонове поле $\bar{\phi}(t)$.
Лінеаризуємо його для малого збурення $\delta\phi = \phi - \bar{\phi}$,
нехтуючи членом $V''(\bar{\phi})\,\delta\phi$ у субхаббловому режимі
($k^2 \gg \mathcal{H}^2$, де масою поля можна знехтувати), і переходимо
до конформного часу $\eta$ та Фур'є-простору. Використовуючи
$\dot{f} = f'/a$ і $\ddot{f} = (f'' - \mathcal{H}f')/a^2$\footnote{%
З означення $d\eta = dt/a$ маємо $\dot{f} = f'/a$.
Диференціюємо ще раз:
$\ddot{f} = \dfrac{d}{dt}\!\left(\dfrac{f'}{a}\right)
= \dfrac{f''}{a^2} - \dfrac{a'}{a^2}\dfrac{f'}{a}
= \dfrac{1}{a^2}(f'' - \mathcal{H}f')$.},
отримуємо для $k$-ї моди \cite{mukhanov_2005}:
\begin{equation}\label{eq:fourier_field}
  \delta\phi_k'' + 2\mathcal{H}\,\delta\phi_k'
  + k^2\,\delta\phi_k = 0,
\end{equation}
де штрих~--- похідна за $\eta$, $k \equiv |\mathbf{k}|$~---
хвильове число у координатах що рухаються з розширенням
(фізична довжина хвилі $\lambda_\mathrm{phys} = a/k$ росте з часом,
тоді як $k$ залишається сталим). Поведінка кожної моди визначається
співвідношенням між $\lambda_\mathrm{phys}$ та горизонтом Хаббла $H^{-1}$:

\begin{itemize}
  \item \textbf{Субхаббловий режим ($k \gg \mathcal{H}$):}
        Для осцилюючої моди $\delta\phi_k' \sim k\,\delta\phi_k$,
        тому хабблівський член за порядком величини:
        $2\mathcal{H}\delta\phi_k' \sim \mathcal{H}k\,\delta\phi_k
        \ll k^2\delta\phi_k$,
        оскільки $\mathcal{H} \ll k$.
        Рівняння~\eqref{eq:fourier_field} зводиться до звичайного
        гармонічного осцилятора:
        \begin{equation*}
          \delta\phi_k'' + k^2\,\delta\phi_k \approx 0.
        \end{equation*}
        Це рівняння вільного поля у пласкому просторі Мінковського~---
        таке саме, як для фотона або будь-якого іншого квантового поля.
        Його квантовий вакуумний стан добре відомий: це стан з мінімальною
        енергією, де кожна мода $k$ знаходиться у нульових коливаннях з
        амплітудою $\sim 1/\sqrt{2k}$ (у одиницях $\hbar = c = 1$).

        У розширюваному просторі де-Сіттера цей стан називається
        \emph{вакуумом Банча--Девіса} \cite{bunch_davies_1978}~---
        це природний вибір початкового стану: на достатньо малих
        масштабах ($k \gg \mathcal{H}$) кривина простору непомітна,
        і поле «не знає», що воно у де-Сіттері. А отже його вакуум
        збігається зі звичайним вакуумом Мінковського~--- тим самим,
        який ми знаємо з КТП у пласкому просторі. Саме цей стан і
        є природною початковою умовою: ми не вигадуємо нічого нового,
        а просто беремо звичайний квантовий вакуум там, де фізика
        ще «не відчула» розширення. З урахуванням
        масштабного фактора $a$ нормування амплітуди:
        \begin{equation*}
          |\delta\phi_k| \sim \frac{1}{a\sqrt{2k}}.
        \end{equation*}
        Множник $1/a$ має простий фізичний зміст: хвильове число
        $k$ визначене у координатах що рухаються разом з розширенням
        (у тих самих $x^i$, де записана метрика FLRW~--- у них
        галактики стоять на місці, а простір між ними розтягується).
        Тому $k$ не змінюється з часом, але фізична довжина хвилі
        $\lambda_\mathrm{phys} = a/k$ розтягується разом з простором.
        Енергія моди «розмазується» по більшому об'єму, і амплітуда
        спадає як $1/a$~--- так само як температура реліктового
        випромінювання спадає як $1/a$ через космологічне
        розтягування фотонів.
  \item \textbf{Суперхаббловий режим ($k \ll \mathcal{H}$):}
        Градієнтний член $k^2\delta\phi_k$ стає нехтовним
        порівняно з $\mathcal{H}^2\delta\phi_k$, і рівняння
        \eqref{eq:fourier_field} перетворюється з осцилятора
        на рівняння загасання:
        \begin{equation*}
          \delta\phi_k'' + 2\mathcal{H}\,\delta\phi_k' \approx 0.
        \end{equation*}
        Член $k^2$ — це «пружина» осцилятора: саме він змушує
        моду коливатися. Коли мода виходить за горизонт,
        $k/\mathcal{H} \to 0$ і пружина «зникає». Фізично:
        два кінці хвилі більше не перебувають у причинному
        контакті через розширення простору~--- хвиля не може
        «підтримувати» себе і замерзає.
        Розв'язок: $\delta\phi_k = \mathrm{const}$ (домінуюча мода).
        Мода «виходить за горизонт» і заморожується
        \cite{hawking_1982}.
\end{itemize}

У суперхаббловому ліміті амплітуда виходить на стаціонарне
значення \cite{guth_pi_1982}:
\begin{equation*}
  |\delta\phi_k| \;\to\; \frac{H}{\sqrt{2k^3}}.
\end{equation*}

Звідси первинний безрозмірний спектр потужності флуктуацій поля
\cite{mukhanov_chibisov_1981}:
\begin{equation}\label{eq:field_power_spec}
  \mathcal{P}_{\delta\phi}(k)
    \equiv \frac{k^3}{2\pi^2}|\delta\phi_k|^2
    = \left(\frac{H}{2\pi}\right)^{\!2}.
\end{equation}

\textit{Зауваження: де-Сіттерівська температура і плоский спектр.}
Результат~\eqref{eq:field_power_spec} має глибший зміст, ніж
просто «амплітуда замороженої моди». Простір де-Сіттера має
горизонт подій~--- так само, як чорна діра. Спостерігач у такому
просторі не має доступу до інформації поза горизонтом, і це
обмеження, за Гіббонсом і Гокінгом \cite{gibbons_hawking_1977},
призводить до теплового випромінювання вакуумних флуктуацій з
характерною температурою
\begin{equation*}
  T_\mathrm{dS} = \frac{H}{2\pi}.
\end{equation*}
Саме ця температура визначає амплітуду заморожених флуктуацій~---
вираз $H/(2\pi)$ у формулі~\eqref{eq:field_power_spec} є не
збігом, а прямим наслідком квантової термодинаміки горизонту.

З цього випливає ключова властивість спектра: оскільки $H \approx
\mathrm{const}$ під час інфляції, амплітуда $\mathcal{P}_{\delta\phi}
\propto H^2$ практично не залежить від $k$~--- тобто флуктуації
всіх масштабів народжуються з однаковою «температурою».
Це і є спектр Харісона--Зельдовича, майже інваріантний за
масштабом. Невелике відхилення від точної сталості~---
спектральний індекс $n_s \neq 1$~--- виникає тому, що $H$
повільно спадає під час slow-roll, і детально розглядається
в наступному підрозділі.

%% --------------------------------------------------------
\section{Збурення кривини і первинний спектр}
%% --------------------------------------------------------

Збурення $\delta\phi = \phi - \bar\phi$ залежить від того, як
проведена координатна сітка: різні системи координат дають різне
значення $\delta\phi$ в одній і тій самій фізичній події.
Це не фізична різниця~--- це артефакт калібрування.
Тому $\delta\phi_k$ не може бути початковою умовою:
ми б передавали від інфляції не фізику, а вибір координат.

Потрібна \emph{калібрувально-інваріантна} величина~---
така, що однакова в будь-якій калібровці.
Її будують як комбінацію збурення поля і збурення метрики,
що взаємно компенсують калібрувальну свободу.
У конформно-ньютонівській калібровці (де метрика має потенціали
$\Phi$ і $\Psi$ зі збуреної метрики FLRW~--- розд.~\ref{sec:metric})
ця комбінація має вигляд:
\begin{equation}\label{eq:R_definition}
  \mathcal{R}_k \equiv \Psi_k - \frac{H}{\dot{\bar\phi}}\,\delta\phi_k.
\end{equation}
\textbf{Фізичний зміст:} $\mathcal{R}_k$~--- це кривина просторових
гіперповерхонь сталої густини енергії.
Уявіть, що ви «підганяєте» координатну сітку так, щоб
$\delta\phi = 0$ скрізь (уніформна калібровка густини)~---
тоді вся фізична неоднорідність «перейде» у кривину простору,
і саме цю кривину вимірює $\mathcal{R}_k$.
Знак мінус у~\eqref{eq:R_definition}: де поле «вище» за фон
($\delta\phi > 0$), там інфлатон ще не докотився до мінімуму,
реогрів відбудеться пізніше, простір встигне розширитися більше~---
і просторова кривина буде від'ємнішою.

Головна перевага $\mathcal{R}_k$~--- вона \emph{зберігається}
на суперхабблових масштабах незалежно від того, що відбувається
з матерією (реогрів, зміна рівняння стану тощо).
Це робить її ідеальною «початковою умовою», що однозначно
передається від інфляції до гарячого Всесвіту.

\textit{Чому $\mathcal{R}_k$ є сталим за горизонтом?}
Рівняння руху для $\mathcal{R}_k$ у Фур'є-просторі має вигляд
\cite{bardeen_1980}:
\begin{equation}\label{eq:R_eom}
  \mathcal{R}_k'' + \frac{2z'}{z}\,\mathcal{R}_k' + k^2\mathcal{R}_k = 0,
  \qquad z \equiv \frac{a\dot{\bar\phi}}{H}.
\end{equation}
У суперхаббловому режимі ($k \ll \mathcal{H}$) член $k^2\mathcal{R}_k$
стає нехтовним. Рівняння зводиться до
$\mathcal{R}_k'' + (2z'/z)\,\mathcal{R}_k' = 0$,
розв'язок якого~--- $\mathcal{R}_k = \mathrm{const}$ (зростаюча мода)
і $\mathcal{R}_k \propto \int d\eta/z^2$ (загасаюча мода).
Загасаюча мода спадає і нею нехтують; зростаюча мода заморожується
на сталому значенні. Ця сталість не залежить від деталей розігріву і
постінфляційної еволюції~--- тому $\mathcal{R}_k$ є надійною
«початковою умовою», що передається від інфляції до гарячого
Всесвіту незмінною \cite{bardeen_1980}.

Тут нам достатньо знати, що в конформно-ньютонівській калібровці
$\mathcal{R}_k$ лінійно пов'язане з $\delta\phi_k$ через швидкість
кочення класичного фону~--- це безпосередньо випливає з означення
$\mathcal{R}$ через $\Phi$ і $\Psi$ при підстановці рівнянь руху:
\begin{equation}\label{eq:curvature_perturbation}
  \mathcal{R}_k = -\frac{H}{\dot{\bar{\phi}}}\,\delta\phi_k.
\end{equation}
Підставляючи~\eqref{eq:field_power_spec}, отримуємо первинний
спектр кривини:
\begin{equation}\label{eq:curvature_power_spec}
  \mathcal{P}_{\mathcal{R}}(k)
    = \left(\frac{H}{\dot{\bar{\phi}}}\right)^{\!2}
      \mathcal{P}_{\delta\phi}
    = \frac{H^4}{4\pi^2\dot{\bar{\phi}}^2}.
\end{equation}

Ця формула є центральною для всього подальшого аналізу. У
стандартному slow-roll режимі $\dot{\bar\phi}$ змінюється
повільно, тому $\mathcal{P}_{\mathcal{R}}(k)$ є майже не
залежним від $k$~--- так званий спектр Харісона--Зельдовича.
Відхилення від плоского спектру параметризується спектральним
індексом:
\begin{equation}\label{eq:ns}
  n_s - 1 \equiv \frac{d\ln\mathcal{P}_{\mathcal{R}}}{d\ln k}
            = -2\varepsilon - \eta_V,
\end{equation}
де $\varepsilon$ та $\eta_V$~--- параметри slow-roll з
формули~\eqref{eq:slow_roll_pot}, обчислені в момент виходу
моди $k$ за горизонт.

\textit{Виведення~\eqref{eq:ns}.}
Кожна мода $k$ виходить за горизонт у момент $k = \mathcal{H}
= aH$, тобто число $e$-folds до кінця інфляції $N_k$ залежить
від $k$: $dN = -H\,dt$, звідки $d\ln k = dN$ вздовж траєкторії
slow-roll. Тому:
\begin{equation*}
  n_s - 1
    = \frac{d\ln\mathcal{P}_{\mathcal{R}}}{d\ln k}
    = \frac{d}{dN}\ln\!\left(\frac{H^4}{\dot{\bar\phi}^2}\right)
    = 4\frac{\dot H}{H^2} - 2\frac{\ddot{\bar\phi}}{H\dot{\bar\phi}}
    = -4\varepsilon + 2\eta_\mathrm{H},
\end{equation*}
де використано $\dot H = -\varepsilon H^2$ (з означення~\eqref{eq:eps})
та $\eta_\mathrm{H} \equiv -\ddot{\bar\phi}/(H\dot{\bar\phi})$
(знак мінус у означенні).
У потенціальному наближенні slow-roll зв'язок між $\eta_\mathrm{H}$
і $\eta_V$ отримується диференціюванням attractor-рівняння
$3H\dot{\bar\phi} \approx -V'(\bar\phi)$ за часом $t$:
\begin{equation*}
  3\dot H\dot{\bar\phi} + 3H\ddot{\bar\phi} \approx -V''(\bar\phi)\dot{\bar\phi}.
\end{equation*}
Ділимо на $3H^2\dot{\bar\phi}$:
\begin{equation*}
  \frac{\dot H}{H^2} + \frac{\ddot{\bar\phi}}{H\dot{\bar\phi}}
  \approx -\frac{V''}{3H^2}
  \approx -\frac{M_\mathrm{Pl}^2 V''}{V} \cdot \frac{V}{3H^2 M_\mathrm{Pl}^2}
  \approx -\eta_V \cdot 1,
\end{equation*}
де в останньому кроці $V/(3H^2 M_\mathrm{Pl}^2) \approx 1$ у slow-roll.
Отже, з $\dot H/H^2 = -\varepsilon$ і $\eta_\mathrm{H} = -\ddot{\bar\phi}/(H\dot{\bar\phi})$:
\begin{equation*}
  -\varepsilon - \eta_\mathrm{H} \approx -\eta_V
  \quad\Rightarrow\quad
  \eta_\mathrm{H} \approx \varepsilon - \eta_V.
\end{equation*}
Точніший рахунок з урахуванням поправок порядку $\varepsilon^2$
дає \cite{mukhanov_2005}:
\[
  \eta_\mathrm{H} \approx \varepsilon - \tfrac{1}{2}\eta_V,
\]
звідки:
\begin{equation*}
  n_s - 1 = -4\varepsilon + 2\bigl(\varepsilon -
  \tfrac{1}{2}\eta_V\bigr) = -2\varepsilon - \eta_V. \qquad\square
\end{equation*}
Planck~2018 вимірює $n_s = 0.9649 \pm 0.0042$~--- ледь нахилений,
але не плоский спектр, що є прямим підтвердженням slow-roll.

Проте якщо на шляху інфлатона є локальна аномалія потенціалу
(сходинка, ділянка USR), швидкість $\dot{\bar\phi}$ різко падає.
Оскільки $\mathcal{P}_{\mathcal{R}} \propto \dot{\bar\phi}^{-2}$,
флуктуації, що виходять за горизонт у цей момент, зазнають
локального підсилення~--- формується \textbf{зламаний спектр}.
Механізм і наслідки детально розглядаються у
§\,\ref{sec:future_research}.

%% --------------------------------------------------------
\section{Заморожування та розморожування мод}
%% --------------------------------------------------------

Еволюція первинних збурень підпорядковується чіткій часовій
схемі:
\begin{enumerate}
  \item \textbf{Заморожування (інфляція):} Простір розширюється
        прискорено ($a \propto e^{Ht}$), горизонт Хаббла
        $H^{-1}$ залишається майже сталим. Масштаби збурень
        послідовно покидають горизонт ($\lambda_\mathrm{phys} >
        H^{-1}$), і збурення $\mathcal{R}_k$ заморожуються,
        зберігаючи відбиток локальної мікрофізики $V(\phi)$
        \cite{bardeen_1980}.
  \item \textbf{Розморожування (постінфляційні епохи):}
        Розширення сповільнюється, горизонт Хаббла росте
        швидше за $a$. Збурення послідовно повертаються в
        горизонт~--- спочатку малі масштаби (великі $k$),
        пізніше великі (малі $k$) \cite{mukhanov_2005}.
\end{enumerate}

Повернувшись в горизонт, кожна мода $\mathcal{R}_k$ стає
початковою умовою для класичної гідродинаміки фотон-баріонної
плазми. Але перш ніж плазма виникне~--- потрібен розігрів:
перетворення енергії інфлатона на гарячі частинки Стандартної
моделі. Саме з цього починається наступний розділ.

%% --------------------------------------------------------
\section{Динаміка первинних гравітаційних хвиль}
%% --------------------------------------------------------

Вияснимо, як еволюціонує тензорні збурення $h_{ij}$. Підставляючи тензорну
частину~\eqref{eq:perturbed_metric} у лінеаризовані рівняння
Айнштайна без джерел анізотропного тиску, отримуємо хвильове
рівняння для кожної Фур'є-моди:
\begin{equation}\label{eq:tensor_wave_eq}
  h_{ij}'' + 2\mathcal{H}\,h_{ij}' + k^2\,h_{ij} = 0.
\end{equation}
Його структура, як ми пізніше побачимо, ідентична рівнянню~\eqref{eq:fourier_field}
для скалярних флуктуацій~--- той самий осцилятор із
хабблівським тертям $2\mathcal{H}$. Тому і результат
аналогічний: моди заморожуються за горизонтом під час
інфляції, і їхній первинний спектр визначається масштабом
Хаббла \cite{starobinsky_1979}:
\begin{equation}\label{eq:tensor_power_spec}
  \mathcal{P}_h(k)
    = \frac{2}{\pi^2 M_\mathrm{Pl}^2}
      \left(\frac{H}{2\pi}\right)^{\!2}
    = \frac{H^2}{\pi^2 M_\mathrm{Pl}^2}.
\end{equation}
Множник $M_\mathrm{Pl}^{-2}$ відрізняє тензорний спектр від
скалярного~\eqref{eq:curvature_power_spec}: гравітаційні хвилі
пригнічені планківською масою, що відображає слабкість
гравітаційної взаємодії.

Відношення амплітуд тензорних і скалярних збурень~---
\emph{тензорно-скалярне відношення}~--- безпосередньо
вимірюється через B-моди поляризації CMB:
\begin{equation}\label{eq:tensor_scalar_ratio}
  r \equiv \frac{\mathcal{P}_h(k_*)}
               {\mathcal{P}_{\mathcal{R}}(k_*)}
    = 16\varepsilon,
\end{equation}
де $k_*$~--- опорний масштаб. Сучасна верхня межа
$r < 0.036$ \cite{bicep_keck_2021} обмежує енергетичний
масштаб інфляції і відсікає моделі з крутим потенціалом.
Зв'язок між $r$ і B-модами детально розглядається у
§\,\ref{sec:pol}.